\section{Wnioski i dalsze kroki}

\subsection{Podsumowanie}

\begin{itemize}[noitemsep]
  \item \textbf{TF-IDF $\rightarrow$ MLP} to sieć neuronowa zbudowana w całości
        od zera (bez fine-/feature-tuningu), zgodna z wymogiem przedmiotu SSN.
        Działa na tych samych splitach, konwencji etykiet i wektoryzatorze co
        baseline, a cała logika jest reużywalna w \texttt{src/model/mlp.py}.
  \item Model osiąga F1 i ROC-AUC na poziomie $\sim$0{,}99 na zbiorze testowym,
        co jest wynikiem porównywalnym z baselinem (regresją logistyczną). Na tym leksykalnie
        łatwym zbiorze sieć nie zyskuje istotnej przewagi nad modelem liniowym.
  \item Diagnoza \textbf{bias-variance} (krzywa uczenia + sweep regularyzacji)
        wskazuje reżim \textbf{niskiego biasu i niskiej wariancji} przy suficie
        metryk zbioru: powiększanie modelu ani zbioru przyniosłoby co najwyżej
        marginalną poprawę, a regularyzacja jedynie domyka niewielką lukę train-val.
  \item Przykładowe predykcje ujawniają, że model jest \textbf{przepewny}:
        prawdopodobieństwa są nasycone na krańcach również przy pomyłkach (słaba
        kalibracja mimo wysokiego ROC-AUC). Nieliczne błędy dotyczą tekstów
        stylistycznie ,,pogranicznych'' (agencyjny styl w artykułach Fake, format
        listy lub tematyka rozrywkowa w artykułach Real).
\end{itemize}

\subsection{Ograniczenia}

Najważniejszym ograniczeniem jest \textbf{sufit metryczny} zbioru: po usunięciu
wycieku klasy nadal różnią się na tyle silnie leksykalnie, że proste reprezentacje
(TF-IDF) wystarczają do niemal perfekcyjnej klasyfikacji. Utrudnia to wykazanie
przewagi bardziej złożonych architektur na tym konkretnym zbiorze. Model
worka słów (TF-IDF) z założenia ignoruje też \textbf{kolejność i kontekst} słów.

\subsection{Dalsze kroki}

Projekt dostarcza zweryfikowaną pętlę treningową PyTorcha (Adam,
\texttt{BCEWithLogitsLoss}, early stopping, dropout, weight decay, obsługa GPU) oraz
kompletny, reużywalny kod modelu w \texttt{src/model/mlp.py}. Naturalnym kierunkiem
dalszego rozwoju, poza zakresem niniejszego projektu, jest architektura
\textbf{BiLSTM + Attention} z osadzeniami \textbf{GloVe~300d}: w odróżnieniu od worka
słów modelowałaby kontekst sekwencyjny, a mechanizm atencji umożliwiłby
\textbf{wyjaśnialność na poziomie pojedynczego artykułu} (wizualizację wag atencji).
